\documentclass[UTF8,zihao=-4]{ctexart}
\usepackage{geometry,xcolor,amsmath,amssymb,enumitem,needspace,setspace,xeCJKfntef,fontspec}
\usepackage{fancyhdr,tabularx,multicol,tikz,hyperref}
\geometry{a4paper,left=18mm,right=16mm,top=16mm,bottom=20mm,footskip=10mm}

% Font profile: preserve the reference font stack where available, but keep
% public/cloud builds reproducible without redistributing proprietary fonts.
\IfFileExists{C:/Windows/Fonts/simsun.ttc}{%
 \newfontfamily\wblatin[Path=C:/Windows/Fonts/]{times.ttf}%
 \newCJKfontfamily\wbsong[Path=C:/Windows/Fonts/]{simsun.ttc}%
 \newCJKfontfamily\wbhei[Path=C:/Windows/Fonts/]{simhei.ttf}%
 \newCJKfontfamily\wbkai[Path=C:/Windows/Fonts/]{simkai.ttf}%
}{%
 \newfontfamily\wblatin{Tinos}%
 \newCJKfontfamily\wbsong{Noto Serif CJK SC}%
 \newCJKfontfamily\wbhei{Noto Sans CJK SC}%
 \newCJKfontfamily\wbkai{AR PL KaitiM GB}%
}
\wblatin

\definecolor{accent}{HTML}{56695F}
\definecolor{muted}{HTML}{626262}
\definecolor{rulegray}{gray}{0.63}
\color{black}
\setstretch{1}
\setlength{\parindent}{0pt}
\setlength{\parskip}{0pt}
\setlength{\emergencystretch}{1.5em}
\setlength{\tabcolsep}{0pt}
\setlength{\columnsep}{8mm}
\setlength{\multicolsep}{0pt}
\AtBeginDocument{\wbsong\fontsize{10.5}{15}\selectfont}
\hypersetup{pdftitle={青云阁教育·四科练习册样张},pdfauthor={青云阁教育}}

\newcommand{\marked}[1]{\CJKunderdot[depth=0.15em,sep=0.01em]{#1}}
\newcommand{\bookornament}{%
 \begin{tikzpicture}[x=1mm,y=1mm,baseline=-.3mm,line cap=round,line join=round]
 \draw[accent,line width=.4pt]
 (0,.4)--(0,4.2)..controls(1.8,5)and(3.5,4.8)..(5,3.7)
 ..controls(6.5,4.8)and(8.2,5)..(10,4.2)--(10,.4);
 \draw[accent,line width=.4pt]
 (0,.4)..controls(1.8,1.2)and(3.5,1)..(5,0)
 ..controls(6.5,1)and(8.2,1.2)..(10,.4);
 \draw[accent,line width=.35pt] (5,0)--(5,3.7);
 \end{tikzpicture}}

% Footer: publication output contains only brand + quotation + page number.
\pagestyle{fancy}\fancyhf{}
\renewcommand{\headrulewidth}{0pt}
\renewcommand{\footrulewidth}{0pt}
\newcommand{\footerquote}{腹有诗书气自华}
\newcommand{\setfooterquote}[1]{\renewcommand{\footerquote}{#1}}
\fancyfoot[L]{%
 {\fontsize{10}{12}\selectfont\wbkai\color{muted}青云阁}%
 \kern2pt{\fontsize{7.8}{10}\selectfont\wbhei\color{muted}教育}%
 \hspace{3mm}\textcolor{rulegray}{\rule[2.5pt]{8mm}{.35pt}}%
 \hspace{3mm}{\fontsize{8}{10}\selectfont\wbsong\color{muted}\footerquote}}
\fancyfoot[R]{\fontsize{9}{12}\selectfont\wblatin\color{muted}\thepage}

\newcounter{exercise}
\newcommand{\titleface}{\wbhei}

% Header keeps the reference geometry while protecting the right-side rail
% from being displaced by long titles.
\newcommand{\pagehead}[3]{%
 \begingroup
 \renewcommand{\arraystretch}{1}%
 \noindent\begin{tabularx}{\linewidth}{@{}X>{\raggedleft\arraybackslash}p{38mm}@{}}
 {\fontsize{9}{12}\selectfont\wbhei\color{muted}#1} & {\fontsize{9}{12}\selectfont\wbsong\color{muted}#3}\\[1.8mm]
 {\fontsize{16}{20}\selectfont\titleface #2} & \bookornament
 \end{tabularx}\par
 \endgroup
 \vspace{2mm}
 \noindent\hbox to\linewidth{\textcolor{accent}{\rule{14mm}{.7pt}}\hspace{3mm}%
 {\color{rulegray}\leaders\hrule height .35pt\hfill}}\par\vspace{3mm}}

\newcommand{\subhead}[1]{\Needspace{3\baselineskip}\par\vspace{3mm}%
 {\textcolor{accent}{\rule[3pt]{3mm}{.6pt}}\hspace{2mm}%
 \fontsize{11}{16}\selectfont\wbhei #1}\par\vspace{2mm}}

% Shared horizontal grid.
% Number x=18 mm; score/options/sub-question x=23 mm.
% Stem stays near x=39.05 mm, preserving the reference stem position.
\newlength{\qnumwidth}
\newlength{\qscorewidth}
\newlength{\qstemgap}
\setlength{\qnumwidth}{5mm}
\setlength{\qscorewidth}{14.85mm}
\setlength{\qstemgap}{1.2mm}

% #1 is retained only for source compatibility and is never printed.
\newcommand{\q}[3]{%
 \Needspace{2.8\baselineskip}%
 \par\stepcounter{exercise}\begingroup
 \normalfont\wbsong\fontsize{10.5}{15}\selectfont
 \parindent=0pt
 \hangindent=\dimexpr\qnumwidth+\qscorewidth+\qstemgap\relax
 \hangafter=1\relax
 \noindent
 \makebox[\qnumwidth][l]{\sffamily\bfseries\arabic{exercise}.}%
 \makebox[\qscorewidth][l]{{\color{muted}（#2分）}}%
 \hspace{\qstemgap}#3\par
 \endgroup\vspace{1mm}}

% Two-column choices. A/C start on the same rail as the score.
\newcommand{\opt}[4]{%
 \Needspace{3.2\baselineskip}%
 \begingroup\renewcommand{\arraystretch}{1.14}
 \noindent\hspace*{\qnumwidth}%
 \begin{tabularx}{\dimexpr\linewidth-\qnumwidth\relax}{@{}>{\raggedright\arraybackslash}X@{\hspace{6mm}}>{\raggedright\arraybackslash}X@{}}
 A. #1&B. #2\\ C. #3&D. #4
 \end{tabularx}\par\endgroup\vspace{3mm}}

% Four-column short choices, also aligned to the score rail.
\newcommand{\optfour}[4]{%
 \Needspace{2.2\baselineskip}%
 \begingroup\renewcommand{\arraystretch}{1.12}
 \noindent\hspace*{\qnumwidth}%
 \begin{tabularx}{\dimexpr\linewidth-\qnumwidth\relax}{@{}XXXX@{}}A. #1&B. #2&C. #3&D. #4\end{tabularx}\par
 \endgroup\vspace{3mm}}

% Long choices. Label starts on score rail; wrapped continuation aligns with option text.
\newcommand{\longoptline}[2]{%
 \par\begingroup
 \hangindent=\dimexpr\qnumwidth+5mm\relax\hangafter=1
 \noindent\hspace*{\qnumwidth}\makebox[5mm][l]{#1.}#2\par
 \endgroup}
\newcommand{\longopts}[4]{%
 \Needspace{5.2\baselineskip}%
 \begingroup
 \longoptline{A}{#1}\vspace{.5mm}%
 \longoptline{B}{#2}\vspace{.5mm}%
 \longoptline{C}{#3}\vspace{.5mm}%
 \longoptline{D}{#4}%
 \endgroup\vspace{3mm}}

% Math/material sub-question. Its marker starts on the score/option rail.
\newcommand{\subq}[2]{%
 \Needspace{2.2\baselineskip}%
 \par\begingroup
 \hangindent=\dimexpr\qnumwidth+7mm\relax\hangafter=1
 \noindent\hspace*{\qnumwidth}\makebox[7mm][l]{（#1）}#2\par
 \endgroup}

\newcommand{\blank}[1]{\rule[-.3ex]{#1}{.4pt}}
\newcount\samplelinecount
\newcommand{\writinglines}[1]{\par\begingroup\offinterlineskip\samplelinecount=0\loop\ifnum\samplelinecount<#1
 \vbox to8mm{\vfill\color{rulegray}\hrule width\linewidth height .4pt}\advance\samplelinecount by1\repeat\endgroup}

\begin{document}

% Chinese
\setcounter{exercise}{0}
\renewcommand{\titleface}{\wbkai}
\setfooterquote{读书百遍，其义自见}
\pagehead{语文\quad 古代诗文阅读}{晋文公与里凫须}{阅读与选择}
\begingroup\wbsong\fontsize{10.5}{16}\selectfont\setlength{\parindent}{2em}
晋文公重耳亡，过曹，里凫须从，因盗重耳资而亡。重耳无粮，馁不能行，子推割股肉以食重耳，然后能行。及重耳反国，国中多不附重耳者。于是里凫须造见，曰：“臣能安晋国。”文公使人应之曰：“子尚何面目来见寡人！欲安晋国也！”里凫须曰：“君沐邪？”使者曰：“否。”凫须曰：“臣闻沐者其心倒，心倒者其言悖。今君不沐，何言之悖也？”使者以闻，文公见之。里凫须仰首曰：“\CJKunderline{离国久，臣民多过君；君反国，而民皆自危。}里凫须又袭竭君之资，避于深山，而君以馁，介子推割股，天下莫不闻。臣之为贼亦大矣，罪至十族，未足塞责；然君诚赦之罪，与骖乘，游于国中，百姓见之，必知不念旧恶，人自安矣。”于是文公大悦，从其计，使骖乘于国中。百姓见之，皆曰：“夫里凫须且不诛而骖乘，吾何惧也？”是以晋国大宁。故书云：“文王卑服，即康功田功。”\par
\endgroup
\hfill{\wbsong\fontsize{9}{12}\selectfont\color{muted}（选自《韩诗外传》）}\par\vspace{3mm}
\q{19}{3}{对下列加点字的注音，不正确的一项是（\quad）。}
\opt{重耳无粮，\marked{馁}不能行（něi）}{子推割股肉以\marked{食}重耳（shí）}{使\marked{骖}乘于国中（cān）}{然君诚\marked{赦}之罪，与骖乘（shè）}
\q{20}{3}{对下列加点词的解释，不正确的一项是（\quad）。}
\opt{晋文公重耳\marked{亡}：亡，逃亡。}{天下\marked{莫}不闻：莫，没有谁。}{于是里凫须\marked{造}见：造，到；去。}{臣之为\marked{贼}亦大矣：贼，小偷。}
\q{21}{3}{下列各组句子中加点词的解说，不正确的一项是（\quad）。}
\longopts{“子推割股肉\marked{以}食重耳”，“以”是连词，表示目的。}{“君沐\marked{邪}”，“邪”是语气词，表示疑问。}{“\marked{然}君诚赦之罪，与骖乘”，“然”是连词，表示转折。}{“夫里凫须\marked{且}不诛而骖乘”，“且”是连词，表示并列。}
\q{23}{3}{下列对文中画横线句子的翻译，正确的一项是（\quad）。}
\longopts{离国太久，臣民大多超过您；您重返国家，百姓都自觉处境危险。}{离国太久，臣民大多怪罪于您；您重返国家，百姓都自觉处境危险。}{离国太久，臣民大多超过您；您反叛国家，百姓都自觉处境危险。}{离国太久，臣民大多怪罪于您；您反叛国家，百姓都自觉处境危险。}
\q{24}{3}{下列对文中有关内容的理解，不正确的一项是（\quad）。}
\longopts{里凫须曾经偷窃流亡中的公子重耳的资财，导致公子重耳断粮。}{公子重耳刚即位为晋国国君，就得到了晋国臣子和百姓的拥戴。}{里凫须求见晋文公，说他可以安定晋国，最初遭到晋文公拒绝。}{晋文公优待里凫须，以此表明自己不念旧恶，安定了晋国民心。}

% Mathematics
\clearpage
\setcounter{exercise}{0}
\renewcommand{\titleface}{\wbhei}
\setfooterquote{纸上得来终觉浅，绝知此事要躬行}
\pagehead{数学\quad 基础与综合}{函数与解析几何}{选择·填空·解答}
\q{2}{8}{函数 $y=\sin 3x\cos 3x$ 是（\quad）。}
\opt{最小正周期为 $\frac{\pi}{3}$ 的奇函数}{最小正周期为 $\frac{\pi}{3}$ 的偶函数}{最小正周期为 $\frac{2\pi}{3}$ 的奇函数}{最小正周期为 $\frac{2\pi}{3}$ 的偶函数}
\vspace{4mm}
\q{10}{8}{函数 $y=-x^2-5x+6$ 的单调递减区间为\blank{40mm}。}
\vspace{10mm}
\subhead{解答题}
\q{14}{18}{已知椭圆 $C$ 的中心为原点，焦点 $F_1,F_2$ 在 $x$ 轴上，离心率为 $\frac57$，点 $(0,2\sqrt6)$ 在 $C$ 上。}
\subq{1}{求 $C$ 的方程。}
\vspace{55mm}
\subq{2}{设点 $M$ 在 $C$ 上，且 $\angle F_1MF_2=90^\circ$，求 $\triangle F_1MF_2$ 的面积。}
\vspace{65mm}

% Politics
\clearpage
\setcounter{exercise}{0}
\setfooterquote{苟日新，日日新，又日新}
\pagehead{政治\quad 基础知识与材料分析}{识记与材料运用}{选择·填空·问答}
\q{2}{3}{在科学社会主义发展史上，无产阶级夺取政权的第一次伟大尝试是（\quad）。}
\opt{十月革命}{成立巴黎公社}{辛亥革命}{建立中华人民共和国}
\q{5}{3}{毛泽东在《〈共产党人〉发刊词》中指出，中国共产党在中国革命中战胜敌人有三个法宝。这里提到的“三个主要的法宝”是（\quad）。}
\opt{武装斗争、工农联盟、党的建设}{统一战线、精兵简政、党的建设}{统一战线、武装斗争、党的建设}{武装斗争、减租减息、党的建设}
\q{29}{3}{2023年10月，在全国宣传思想文化工作会议上，正式提出和系统阐述习近平文化思想，这标志着我们党对\blank{35mm}文化建设规律的认识达到了新高度。}
\subhead{材料问答}
\q{32}{20}{阅读材料，回答问题。}
\begingroup\wbsong\fontsize{10.5}{16}\selectfont\setlength{\parindent}{2em}
我国有49.2万个村民委员会、11.6万个居民委员会，覆盖城乡全体居民。这些居民委员会、村民委员会都是居民群众自己选举产生的。党的二十大报告指出，要积极发展基层民主，健全基层党组织领导的基层群众自治机制，完善基层直接民主制度和工作体系，增强城乡社区群众自我管理、自我服务、自我教育、自我监督的实效。\par\endgroup
\vspace{3mm}
结合材料，运用《政治与法治》知识说明我国基层群众自治及其在基层民主实践中的意义。\par
\writinglines{10}

% English
\clearpage
\setcounter{exercise}{0}
\renewcommand{\titleface}{\wblatin\itshape}
\setfooterquote{读书破万卷，下笔如有神}
\pagehead{英语\quad 阅读理解与词汇}{Tickets on Sale for Science Friday}{阅读·单词拼写}
\begingroup\wblatin\fontsize{10.5}{15}\selectfont
\begin{multicols}{2}
\raggedright\hyphenpenalty=10000\exhyphenpenalty=10000
\setlength{\parskip}{2mm}\setlength{\parindent}{0pt}
On April 21, radio and television journalist Ira Flatow will host his public radio show ``Science Friday'' live on campus. The program is broadcast weekly on Public Radio International to an audience of 1.8 million people and, according to the show's website, offers listeners ``a lively, informative discussion on science, technology, health, space and the environment.''

In addition to Flatow as host, the Science Friday event at Miami will feature a special group of interview guests, to be announced in the coming weeks.

The most recent program discusses the physics behind an ice skater's perfect spin and highlights how artificial intelligence is making it hard to tell real news from fake.

\columnbreak
For his live show at Miami University, Flatow will focus on science news and stories local to southwest Ohio. Flatow has been sharing science with the public for over 35 years. Before he started hosting Science Friday in 1991, Flatow was the science journalist for National Public Radio from 1971 to 1986. He has reported from the South Pole, Kennedy Space Center, Three Mile Island, Antarctica and, soon, Oxford, Ohio.

Professors are encouraged to involve the Science Friday event in their courses, and students attending the courses will receive discounted \$5 tickets. Tickets are currently on sale to all students for \$10 and to adults for \$20. They can be bought online or in-person at the Box Office in the Campus Avenue Building.
\end{multicols}\endgroup
\vspace{4mm}
\q{31}{4}{What does Flatow do?}
\opt{He's a scientist.}{He's a journalist.}{He's a professor.}{He's a TV actor.}
\q{32}{4}{What will be discussed on the live show at Miami University?}
\opt{The physics behind ice skating.}{Applications of artificial intelligence.}{Science news of southwest Ohio.}{The climate change at the South Pole.}
\q{33}{4}{What is the price of a ticket to the event for adults?}
\optfour{\$5}{\$10}{\$15}{\$20}
\subhead{单词拼写}
\q{46}{2}{After \blank{25mm}（早餐）we went for a walk on the beach.}
\vspace{1mm}
\q{47}{2}{The street is not \blank{25mm}（安全的）for children to play in.}
\vspace{1mm}
\q{48}{2}{He has to drive a long \blank{25mm}（距离）to work every day.}
\vspace{1mm}
\q{49}{2}{You cannot \blank{25mm}（判断）a book by its cover.}
\vspace{1mm}
\q{50}{2}{People worldwide \blank{25mm}（庆祝）the New Year in different ways.}

\end{document}
